\documentclass[10pt,a4paper]{amsart}
\usepackage[margin=19mm]{geometry}
\usepackage{amsmath,amssymb,mathtools}
\usepackage[scr=rsfs,cal=euler]{mathalfa}
\usepackage{xfrac}
\usepackage[unicode,hidelinks,bookmarks=false]{hyperref}
\newcommand{\ie}{i.e.\ }
\newcommand{\cf}{cf.\ }
\newcommand{\resp}{respectively\ }
\newcommand{\Subsection}{\subsection}
\providecommand{\nobreakdash}{\nobreak\hbox{-}}
\setlength{\emergencystretch}{2em}
\multlinegap0pt
\usepackage{amssymb}
\usepackage{xcolor}
\usepackage{algorithm}\usepackage{algpseudocode}
\newtheorem{theorem}{Theorem}
\newcommand{\N}{\mathcal{N}}
\newcommand{\bp}{\mathbf{p}}
\newcommand{\cP}{\mathcal{P}}
\title{An Information Theoretic Treatment of Yager's Probability Distribution Negation}
\author{Roberto Bruno, Ugo Vaccaro}
\date{Source: arXiv:2608.00594v1 (CC BY 4.0)}
\begin{document}
\maketitle

\noindent\emph{Source and licence.} An Information Theoretic Treatment of Yager's Probability Distribution Negation. Version arXiv:2608.00594v1 (CC BY 4.0), distributed by arXiv under the Creative Commons Attribution 4.0 International licence, http://creativecommons.org/licenses/by/4.0/, as recorded in the arXiv licence metadata for this exact version and shown on the abstract page https://arxiv.org/abs/2608.00594v1. Full licence notice: \texttt{licenses/arxiv-2608.00594-CC-BY-4.0.txt}.\par\smallskip

\noindent\textit{Editorial scope.} Counted unit: the article's theorem th:majorization (source0.tex lines 516--522). The article's complete original proof of it is delivered with it (source0.tex lines 523--554, one \texttt{proof} environment of 32 lines). No mathematics is written, repaired or re-derived: the statement and the proof are the article's own lines, in the article's own notation, and no retained line is substituted or re-flowed.  Retained uncounted as the source-local context the proof needs: lines 308--312.  Nothing was omitted from any retained span, so the package carries no omission record.  No retained line contains a citation, so the wrapper carries no bibliography and no bibliography entry.  No retained line contains a figure environment, an image-inclusion command or drawing code, so no image is delivered and the package declares no asset dependency.  Editorial formatting only, in addition: an AMS \texttt{amsart} preamble that restates the article's own declarations and macros used by the retained lines, namely the environments \texttt{theorem} on separate counters, together with the standard packages those lines need.  The retained environments print in this document as Theorem 1 in order of appearance, whereas the article prints the same items with the numbering of its own full document.  The complete native source of the article as distributed in the arXiv e-print for this exact version is bundled unchanged as \texttt{sources/206552/source0.tex}.
\begin{align}
    &0 \leq N(p_i) \leq 1, \forall i=1,\dots,n,\\
    &\sum_{i=1}^n N(p_i)=1,\\
    &\text{if } p_i\leq p_j, \text{ then } N(p_i) \geq N(p_j). \label{mon}
\end{align}

\begin{theorem}\label{th:majorization}
Let $\bp\in\cP_n$ and let $N_\alpha$ and $N_{\alpha'}$, with $\alpha, \alpha'\in[0,1]$, with 
$\alpha' <\alpha$, be two arbitrary linear negators. Then, $\N_\alpha(\bp)$ is majorized by $\N_{\alpha'}(\bp)$, i.e.,
\begin{equation}\label{majalpha}
\N_\alpha(\bp)\preceq \N_{\alpha'}(\bp).
\end{equation}
\end{theorem}

\begin{proof}
Assume without loss of generality that $\bp$ is ordered in a non-increasing fashion, that is, $p_1\geq\dots\geq p_n$. Let $\alpha'=\alpha-\epsilon$.
The two probability distributions that we obtain from the application of the two negators
 $N_\alpha$ and $N_{\alpha'}$, are the following:

\begin{align*}
     \N_\alpha(\bp) = \Biggl(&\alpha\frac{1}{n} + (1-\alpha)\frac{1-p_1}{n-1},\dots,\\&\ \alpha\frac{1}{n} + (1-\alpha)\frac{1-p_n}{n-1}\Biggr)
\end{align*}
and
\begin{align*}
    \N_{\alpha'}(\bp) = \Biggl(&(\alpha-\epsilon)\frac{1}{n} + (1-\alpha+\epsilon)\frac{1-p_1}{n-1},\dots,\\&(\alpha-\epsilon)\frac{1}{n} + (1-\alpha + \epsilon)\frac{1-p_n}{n-1}\Biggr)
\end{align*}

Since $\bp$ is ordered in a non-increasing fashion, by (\ref{mon}) the probability distributions 
$\N_\alpha(\bp)$ and $\N_{\alpha'}(\bp)$ are ordered in a non-decreasing fashion. Therefore, in order to prove that $\N_\alpha(\bp)\preceq \N_{\alpha'}(\bp)$ we have to show that, for each $k=1,\dots,n$, it holds that
\begin{equation}
    \sum_{i=1}^k N_{\alpha}(p_{n-i+1}) \leq \sum_{i=1}^kN_{\alpha'}(p_{n-i+1}).\label{eq:to_prove}
\end{equation}
For an arbitrary $k\in\{1,\dots,n\}$, it holds that
\begin{align}
    \sum_{i=1}^kN_{\alpha'}(p_{n-i+1}) &= \sum_{i=1}^k \frac{\alpha - \epsilon}{n} \nonumber\\&\quad+ (1-\alpha+\epsilon)\frac{1-p_{n-i+1}}{n-1}\nonumber\\
    &=\sum_{i=1}^k \frac{\alpha}{n} + (1-\alpha)\frac{1-p_{n-i+1}}{n-1}\nonumber\\&\quad+ \epsilon\left(\frac{1-p_{n-i+1}}{n-1} - \frac{1}{n}\right)\nonumber\\
    &=\sum_{i=1}^k \frac{\alpha}{n} + (1-\alpha)\frac{1-p_{n-i+1}}{n-1}\nonumber\\&\quad+ \sum_{i=1}^k \epsilon\left(\frac{1-p_{n-i+1}}{n-1} - \frac{1}{n}\right)\nonumber\\
    &\geq \sum_{i=1}^k \frac{\alpha}{n} + (1-\alpha)\frac{1-p_{n-i+1}}{n-1}\label{in}\\
    &=\sum_{i=1}^k N_{\alpha}(p_{n-i+1}),\label{eq:arbitrary_k}
\end{align}
where the inequality (\ref{in}) follows from the fact that 
$$\sum_{i=1}^k \epsilon\left(\frac{1-p_{n-i+1}}{n-1} - \frac{1}{n}\right)\geq 0$$ since 
$$\left(\frac{1}{n},\dots,\frac{1}{n}\right) \preceq \left(\frac{1-p_n}{n-1},\dots,\frac{1-p_1}{n-1}\right).$$

Thereby, since (\ref{eq:arbitrary_k}) holds for any $k\in\{1,\dots,n\}$, we have that $\N_\alpha(\bp)\preceq\N_{\alpha'}(\bp)$.
\end{proof}

\end{document}
